\section*{Chapitre \ref{ch:g8:combustion-and-fuels} --- La combustion~: combustibles, produits et gaz à effet de serre}

\begin{solution}{exo:g8:combustion-and-fuels:1}
\ce{C + O2 -> CO2}.
\end{solution}

\begin{solution}{exo:g8:combustion-and-fuels:2}
Avec assez de dioxygène, tout le carbone devient du dioxyde de carbone
(\omterm{def:g8:combustion-and-fuels:complete}{combustion complète}). Avec trop peu, une partie devient du monoxyde de
carbone ou de la suie (\omterm{def:g8:combustion-and-fuels:complete}{combustion incomplète}).
\end{solution}

\begin{solution}{exo:g8:combustion-and-fuels:3}
Il n'a ni couleur, ni odeur, ni goût~: personne ne le remarque, et
pourtant il est toxique.
\end{solution}

\begin{solution}{exo:g8:combustion-and-fuels:4}
Fossiles~: le charbon, le gaz naturel, l'essence. Renouvelables~: le
bois, le biogaz, l'éthanol de betterave.
\end{solution}

\begin{solution}{exo:g8:combustion-and-fuels:5}
\ce{C3H8 + 5O2 -> 3CO2 + 4H2O}.
\end{solution}

\begin{solution}{exo:g8:combustion-and-fuels:6}
Carbone~: 2 \ce{CO2}~; hydrogène~: 3 \ce{H2O}~; oxygène à droite $4 + 3 =
7$, dont 1 vient de l'éthanol, donc 3 \ce{O2}~:
\ce{C2H6O + 3O2 -> 2CO2 + 3H2O}.
\end{solution}

\begin{solution}{exo:g8:combustion-and-fuels:7}
De la suie, c'est-à-dire du carbone. La \omterm{def:g4:what-a-fire-needs:burning}{combustion} est incomplète~: la
flamme manque de dioxygène.
\end{solution}

\begin{solution}{exo:g8:combustion-and-fuels:8}
Environ 339 ppm en 1980 et environ 370 ppm en 2000~: une hausse
d'environ 30 ppm.
\end{solution}

\begin{solution}{exo:g8:combustion-and-fuels:9}
Le carbone du bois a été pris à l'\omterm{def:g4:air-a-mixture-of-gases:air}{air} par l'arbre il y a quelques années
ou quelques décennies~: le \omterm{def:g4:what-a-fire-needs:burning}{brûler} rend ce dioxyde de carbone. Le carbone
du charbon était stocké sous terre depuis des millions d'années~: le
\omterm{def:g4:what-a-fire-needs:burning}{brûler} ajoute du dioxyde de carbone nouveau à l'\omterm{def:g4:air-a-mixture-of-gases:air}{air}.
\end{solution}

\begin{solution}{exo:g8:combustion-and-fuels:10}
\ce{2C + O2 -> 2CO}.
\end{solution}

\begin{solution}{exo:g8:combustion-and-fuels:11}
$(427 - 316)/316 = 111/316 \approx 0.35$~: environ \qty{35}{\percent}.
\end{solution}

\begin{solution}{exo:g8:combustion-and-fuels:12}
Au début, il y a assez de dioxygène~: \ce{CH4 + 2O2 -> CO2 + 2H2O}.
À mesure que le dioxygène de la pièce fermée est consommé, la flamme
manque d'\omterm{def:g4:air-a-mixture-of-gases:air}{air}~: \ce{2CH4 + 3O2 -> 2CO + 4H2O}. Le monoxyde de carbone,
inodore et toxique, s'accumule dans la pièce.
\end{solution}

\begin{solution}{pb:g8:combustion-and-fuels:1}
\textbf{1.} \ce{C3H8 + 5O2 -> 3CO2 + 4H2O}.

\textbf{2.} À gauche~: 3 C, 8 H, 10 O. À droite~: 3 C, $4 \times 2 = 8$ H,
$3 \times 2 + 4 = 10$ O. L'équation est ajustée.

\textbf{3.} \ce{2C3H8 + 7O2 -> 6CO + 8H2O}.

\textbf{4.} La \omterm{def:g8:combustion-and-fuels:complete}{combustion complète}~: 5 \omterm{def:g7:atoms-and-molecules:molecule}{molécules} de dioxygène par
\omterm{def:g7:atoms-and-molecules:molecule}{molécule} de propane, contre $7/2 = 3.5$ pour la \omterm{def:g8:combustion-and-fuels:complete}{combustion incomplète}.

\textbf{5.} L'\omterm{def:g4:air-a-mixture-of-gases:air}{air} de la tente fermée manque bientôt de dioxygène~; la
flamme n'en reçoit plus assez et \omterm{def:g4:what-a-fire-needs:burning}{brûle} de façon incomplète, en jaune.

\textbf{6.} Le monoxyde de carbone~: il n'a pas de couleur, pas d'odeur
(et pas de goût).

\textbf{7.} GHS02, la flamme~: extrêmement inflammable. GHS04, la
bouteille de gaz~: un gaz conservé sous pression.

\textbf{8.} Par exemple~: ne l'utiliser qu'en plein \omterm{def:g4:air-a-mixture-of-gases:air}{air}, jamais dans une
tente ni dans une pièce fermée~; le tenir éloigné de tout ce qui peut
\omterm{def:g4:what-a-fire-needs:burning}{brûler}~; changer la cartouche loin de toute flamme.

\textbf{9.} $132 \div 44 = 3$ fois plus lourd.

\textbf{10.} Du dioxygène de l'\omterm{def:g4:air-a-mixture-of-gases:air}{air}, qui se retrouve dans le dioxyde de
carbone (et dans l'eau).

\textbf{11.} Un \omterm{def:g8:combustion-and-fuels:fossil}{combustible fossile} (il provient du gaz naturel ou du
pétrole)~: le \omterm{def:g4:what-a-fire-needs:burning}{brûler} ajoute du dioxyde de carbone à l'\omterm{def:g4:air-a-mixture-of-gases:air}{air}.

\textbf{12.} $450 \times 3 = \qty{1350}{g} = \qty{1.35}{kg}$ de dioxyde
de carbone.
\end{solution}
